\section{Methodology}

A localized perturbation in the free surface height was initialized as a circular bump centered in the domain. The initial condition is a localized Gaussian perturbation centered in the domain:
\begin{equation}
h(x,y,0) = 10 \exp\left(-\frac{(x-x_c)^2 + (y-y_c)^2}{0.01\,L_x^2}\right),
\end{equation}
where $x_c$ and $y_c$ was chosen to be at the domain center. 

Solid wall boundary conditions were applied at all domain boundaries. This implies that the velocity components vanish at the boundaries as 
\begin{equation}
u = 0 \quad \text{at } x=0,\,L_x, \qquad
v = 0 \quad \text{at } y=0,\,L_y,
\end{equation}
which leads to reflection of waves at the domain edges.

The computational domain was discretized using a uniform grid with the following parameters:
\begin{align}
N_x &= 100, & N_y &= 100, \\
L_x &= 100000\,\text{m}, & L_y &= 100000\,\text{m}, \\
L_t &= 1000\,\text{s}.
\end{align}
Other useful parameters were set to $f = 10^{-4}\ \mathrm{s^{-1}}$, $g = 10\ \mathrm{m\,s^{-2}}$, $H_0 = 1000\ \mathrm{m}$, $\mu = 0.45$, and $\gamma = 0.05$.

Model output was saved every $20$ time steps in order to reduce storage requirements while still resolving the temporal evolution of the solution.
The total energy and volume were computed and stored at every time step as
\begin{equation}
V = \sum_{i=1}^{N_x} \sum_{j=1}^{N_y} h_{i,j} \, \Delta x \, \Delta y,
\end{equation}

\begin{equation}
E_P = \frac{1}{2} g \sum_{i=1}^{N_x} \sum_{j=1}^{N_y} h_{i,j}^2 \, \Delta x \, \Delta y,
\end{equation}

\begin{equation}
E_K = \frac{1}{2} H_0 \sum_{i=1}^{N_x} \sum_{j=1}^{N_y} \left( u_{i,j}^2 + v_{i,j}^2 \right) \Delta x \, \Delta y,
\end{equation}

where $V$ is the total volume, $E_P$ is the potential energy, and $E_K$ is the kinetic energy of the system.

The Fortran code used to simulate this scheme can be seen in \autoref{code}. 